\documentclass[10pt]{article}
\usepackage[utf8]{inputenc}
\usepackage[T1]{fontenc}
\usepackage{amsmath}
\usepackage{amsfonts}
\usepackage{amssymb}
\usepackage[version=4]{mhchem}
\usepackage{stmaryrd}

\title{ECN 711: More on Taxes and Debt. Production }

\author{Gustavo Ventura}
\date{}


\begin{document}
\maketitle
Taxes and Debt - Remarks

\section*{Taxes and Debt - Remarks}
\begin{itemize}
  \item Consider two tax policies, that differ in the taxes assigned to each individual at different dates. These policies, however satisfy:
\end{itemize}

$$
\tau_{t}^{h}(t)+\frac{\tau_{t}^{h}(t+1)}{R(t)}=\hat{\tau}_{t}^{h}(t)+\frac{\hat{\tau}_{t}^{h}(t+1)}{R(t)},
$$

all $h$ and all $t \geq 1$.

\section*{Taxes and Debt - Remarks}
\begin{itemize}
  \item Consider two tax policies, that differ in the taxes assigned to each individual at different dates. These policies, however satisfy:
\end{itemize}

$$
\tau_{t}^{h}(t)+\frac{\tau_{t}^{h}(t+1)}{R(t)}=\hat{\tau}_{t}^{h}(t)+\frac{\hat{\tau}_{t}^{h}(t+1)}{R(t)},
$$

all $h$ and all $t \geq 1$.

\begin{itemize}
  \item What is the implication of this?
\end{itemize}

\section*{Taxes and Debt - Remarks}
\begin{itemize}
  \item Consider two tax policies, that differ in the taxes assigned to each individual at different dates. These policies, however satisfy:
\end{itemize}

$$
\tau_{t}^{h}(t)+\frac{\tau_{t}^{h}(t+1)}{R(t)}=\hat{\tau}_{t}^{h}(t)+\frac{\hat{\tau}_{t}^{h}(t+1)}{R(t)},
$$

all $h$ and all $t \geq 1$.

\begin{itemize}
  \item What is the implication of this?
  \item Clearly, for a given path of interest rates, individual choices will be the same in the two scenarios. Why?
\end{itemize}

Taxes and Debt - Remarks

\section*{Taxes and Debt - Remarks}
\begin{itemize}
  \item This is the basis of Ricardian Equivalence. Policies that keep the present value of individual tax liabilities the same are equivalent.
\end{itemize}

\section*{Taxes and Debt - Remarks}
\begin{itemize}
  \item This is the basis of Ricardian Equivalence. Policies that keep the present value of individual tax liabilities the same are equivalent.
  \item The name comes from David Ricardo (1772-1823) who first noted this and after discussing the modes to pay for war, noted:\\
''In point of economy, there is no real difference in either of the modes ... A man who has 10000 pounds, paying him an income of 500 pounds, out of which he has to pay 100 pounds per annum towards the interest of the debt, is really worth only 8,000 pounds , and would be equally rich, whether he continued to pay 100 pounds per annum or at once, and only once, sacrificed 2,000."
\end{itemize}

Taxes and Debt - Remarks

\section*{Taxes and Debt - Remarks}
\begin{itemize}
  \item Put differently, changes in the timing of tax collections under key assumptions is irrelevant in a particular sense, or neutral.
\end{itemize}

\section*{Taxes and Debt - Remarks}
\begin{itemize}
  \item Put differently, changes in the timing of tax collections under key assumptions is irrelevant in a particular sense, or neutral.
  \item In the context of our economy, one can prove the following result:
\end{itemize}

\section*{Proposition: Ricardian equivalence.}
Fix an "initial" government (tax and debt) policy and consider an associated "initial" competitive equilibrium. Now consider a new tax policy that preserves the initial policy's present value (at the initial equilibrium interest rates) of each individual's tax liability. This policy is equivalent to the initial one in the following sense. There exists a debt policy such that the initial allocation and interest rates, together with the new tax and debt policy, constitute an equilibrium.

\section*{Taxes and Debt - Remarks}
To prove this claim, we need to show:

\section*{Taxes and Debt - Remarks}
To prove this claim, we need to show:\\
(1) Show that initial consumption allocation is still optimal under the initial interest rate sequence and the new taxes.

\section*{Taxes and Debt - Remarks}
To prove this claim, we need to show:\\
(1) Show that initial consumption allocation is still optimal under the initial interest rate sequence and the new taxes.

Easy as we discussed before. If present values of tax liabilities do not change, budget constraints are unchanged under the initial interest rate sequence.

\section*{Taxes and Debt - Remarks}
To prove this claim, we need to show:\\
(1) Show that initial consumption allocation is still optimal under the initial interest rate sequence and the new taxes.

Easy as we discussed before. If present values of tax liabilities do not change, budget constraints are unchanged under the initial interest rate sequence.\\
(2) Show that there exist bond issues such that markets clear under the initial interest rate and the new taxes.

\section*{Taxes and Debt - Remarks}
To prove this claim, we need to show:\\
(1) Show that initial consumption allocation is still optimal under the initial interest rate sequence and the new taxes.

Easy as we discussed before. If present values of tax liabilities do not change, budget constraints are unchanged under the initial interest rate sequence.\\
(2) Show that there exist bond issues such that markets clear under the initial interest rate and the new taxes.

More involved... induction argument.

Production and labor supply

\section*{Production and labor supply}
\begin{itemize}
  \item We now extend the model to allow for production and an endogenous choice of hours of work. There is a single, perishable good in the economy in each date as before.
\end{itemize}

\section*{Production and labor supply}
\begin{itemize}
  \item We now extend the model to allow for production and an endogenous choice of hours of work. There is a single, perishable good in the economy in each date as before.
  \item Preferences: individuals now value streams of consumption (c) and leisure ( $z$ ).
\end{itemize}


\begin{equation*}
u^{h}(.)=v^{h}\left(c^{h}(t), z^{h}(t)\right)+\beta v^{h}\left(c^{h}(t+1), z^{h}(t+1)\right) \tag{1}
\end{equation*}


$v$ is continuous, strictly increasing and strictly concave. $\beta>0$.

Production and labor supply

\section*{Production and labor supply}
\begin{itemize}
  \item Constraints on time use:
\end{itemize}

$$
1=z^{h}(t)+h^{h}(t) ; 1=z^{h}(t+1)+h^{h}(t+1)
$$

where $h^{h}(t)$ stands for labor time of individual $h$ at $t$.

\section*{Production and labor supply}
\begin{itemize}
  \item Constraints on time use:
\end{itemize}

$$
1=z^{h}(t)+h^{h}(t) ; 1=z^{h}(t+1)+h^{h}(t+1)
$$

where $h^{h}(t)$ stands for labor time of individual $h$ at $t$.

\begin{itemize}
  \item Prices for labor services: $w(t)$ and $w(t+1)$.
\end{itemize}

\section*{Production and labor supply}
\begin{itemize}
  \item Constraints on time use:
\end{itemize}

$$
1=z^{h}(t)+h^{h}(t) ; 1=z^{h}(t+1)+h^{h}(t+1)
$$

where $h^{h}(t)$ stands for labor time of individual $h$ at $t$.

\begin{itemize}
  \item Prices for labor services: $w(t)$ and $w(t+1)$.
  \item Individuals maximize ( 1 ), subject to time constraints and lifetime budget constraint:
\end{itemize}

$$
c_{t}^{h}(t)+\frac{c_{t}^{h}(t+1)}{R(t)}=w(t) h^{h}(t)+\frac{w(t+1) h^{h}(t+1)}{R(t)} .
$$

Production and labor supply

\section*{Production and labor supply}
\begin{itemize}
  \item Competitive firms have access to production technology: $y(t)=A(t) h(t)$
\end{itemize}

\section*{Production and labor supply}
\begin{itemize}
  \item Competitive firms have access to production technology: $y(t)=A(t) h(t)$
  \item Note that aggregate output will be:
\end{itemize}

$$
Y(t)=A(t)\left(\sum_{h=1}^{N(t)} h_{t}^{h}(t)+\sum_{h=1}^{N(t-1)} h_{t-1}^{h}(t)\right)
$$

\section*{Production and labor supply}
\begin{itemize}
  \item Competitive firms have access to production technology: $y(t)=A(t) h(t)$
  \item Note that aggregate output will be:
\end{itemize}

$$
Y(t)=A(t)\left(\sum_{h=1}^{N(t)} h_{t}^{h}(t)+\sum_{h=1}^{N(t-1)} h_{t-1}^{h}(t)\right)
$$

\begin{itemize}
  \item Profit maximization will imply that $w(t)=A(t)$.
\end{itemize}

First-order conditions

\section*{First-order conditions}
\begin{itemize}
  \item FOCs for individual maximization imply:
\end{itemize}

$$
\begin{aligned}
w(t) v_{1}(t) & =v_{2}(t) \\
w(t+1) v_{1}(t+1) & =v_{2}(t+1) \\
v_{1}(t) & =\beta R(t) v_{1}(t+1)
\end{aligned}
$$

\section*{First-order conditions}
\begin{itemize}
  \item FOCs for individual maximization imply:
\end{itemize}

$$
\begin{aligned}
w(t) v_{1}(t) & =v_{2}(t) \\
w(t+1) v_{1}(t+1) & =v_{2}(t+1) \\
v_{1}(t) & =\beta R(t) v_{1}(t+1)
\end{aligned}
$$

\begin{itemize}
  \item Two intratemporal conditions, and one intertemporal condition.
\end{itemize}

\section*{First-order conditions}
\begin{itemize}
  \item FOCs for individual maximization imply:
\end{itemize}

$$
\begin{aligned}
w(t) v_{1}(t) & =v_{2}(t) \\
w(t+1) v_{1}(t+1) & =v_{2}(t+1) \\
v_{1}(t) & =\beta R(t) v_{1}(t+1)
\end{aligned}
$$

\begin{itemize}
  \item Two intratemporal conditions, and one intertemporal condition.
  \item What is the interpretation of individual optimality conditions?
\end{itemize}

\section*{Equilibrium}
\section*{Equilibrium}
\begin{itemize}
  \item In equilibrium, $\tilde{w}(t)=A(t)$, all $t$. Likewise,
\end{itemize}

$$
\sum_{h=1}^{N(t)} \tilde{I}^{h}(t)=0
$$

for all $t$.

\section*{Equilibrium}
\begin{itemize}
  \item In equilibrium, $\tilde{w}(t)=A(t)$, all $t$. Likewise,
\end{itemize}

$$
\sum_{h=1}^{N(t)} \tilde{I}^{h}(t)=0
$$

for all $t$.

\begin{itemize}
  \item Exercise: Define a competitive equilibrium. What objects are involved?
\end{itemize}

\section*{Equilibrium}
\begin{itemize}
  \item In equilibrium, $\tilde{w}(t)=A(t)$, all $t$. Likewise,
\end{itemize}

$$
\sum_{h=1}^{N(t)} \tilde{I}^{h}(t)=0
$$

for all $t$.

\begin{itemize}
  \item Exercise: Define a competitive equilibrium. What objects are involved?
  \item Note that aggregate feasibility is as usual, $C(t)=Y(t)$.
\end{itemize}

\section*{Equilibrium}
\begin{itemize}
  \item In equilibrium, $\tilde{w}(t)=A(t)$, all $t$. Likewise,
\end{itemize}

$$
\sum_{h=1}^{N(t)} \tilde{I}^{h}(t)=0
$$

for all $t$.

\begin{itemize}
  \item Exercise: Define a competitive equilibrium. What objects are involved?
  \item Note that aggregate feasibility is as usual, $C(t)=Y(t)$.
  \item Exercise: show that an equilibrium consumption allocation and an equilibrium allocation of labor are feasible if and only if aggregate borrowing and lending is zero at all dates.
\end{itemize}

\section*{Example}
\section*{Example}
\begin{itemize}
  \item Identical individuals within a cohort. Preferences:
\end{itemize}

$$
v(c, z)=\log (c)+\frac{z^{1-\sigma}}{1-\sigma}
$$

$$
\sigma>0, \sigma \neq 1 .
$$

\section*{Example}
\begin{itemize}
  \item Identical individuals within a cohort. Preferences:
\end{itemize}

$$
v(c, z)=\log (c)+\frac{z^{1-\sigma}}{1-\sigma}
$$

$$
\sigma>0, \sigma \neq 1
$$

\begin{itemize}
  \item $A(t+1) / A(t)=(1+g)$.
\end{itemize}

\section*{Example}
\begin{itemize}
  \item Identical individuals within a cohort. Preferences:
\end{itemize}

$$
v(c, z)=\log (c)+\frac{z^{1-\sigma}}{1-\sigma}
$$

$$
\sigma>0, \sigma \neq 1
$$

\begin{itemize}
  \item $A(t+1) / A(t)=(1+g)$.
  \item What is the equilibrium path of interest rates, consumption and hours?
\end{itemize}

\section*{Example}
\section*{Example}
\begin{itemize}
  \item FOC's imply:
\end{itemize}

$$
\begin{aligned}
\frac{w(t)}{c(t)} & =(1-h(t))^{-\sigma} \\
\frac{w(t+1)}{c(t+1)} & =(1-h(t+1))^{-\sigma} \\
\frac{1}{c(t)} & =\beta R(t) \frac{1}{c(t+1)}
\end{aligned}
$$

\section*{Example}
\section*{Example}
\begin{itemize}
  \item In equilibrium, there is no intertemporal trade. Hence,
\end{itemize}

$$
\begin{gathered}
\tilde{c}(t)=\tilde{w}(t) \tilde{h}(t)=A(t) \tilde{h}(t) \\
\tilde{c}(t+1)=\tilde{w}(t+1) \tilde{h}(t+1)=A(t+1) \tilde{h}(t+1)
\end{gathered}
$$

\section*{Example}
\begin{itemize}
  \item In equilibrium, there is no intertemporal trade. Hence,
\end{itemize}

$$
\begin{gathered}
\tilde{c}(t)=\tilde{w}(t) \tilde{h}(t)=A(t) \tilde{h}(t) \\
\tilde{c}(t+1)=\tilde{w}(t+1) \tilde{h}(t+1)=A(t+1) \tilde{h}(t+1)
\end{gathered}
$$

\begin{itemize}
  \item Hours of work are constant over time. $\exists \tilde{h} \in(0,1)$, unique, that solves:
\end{itemize}

$$
\frac{1}{\tilde{h}}=(1-\tilde{h})^{-\sigma}
$$

\section*{Example}
\section*{Example}
\begin{itemize}
  \item Output grows at rate $1+g$.
\end{itemize}

\section*{Example}
\begin{itemize}
  \item Output grows at rate $1+g$.
  \item The interest rate sequence is constant over time:
\end{itemize}

$$
\tilde{R}(t)=\tilde{R}=\frac{1+g}{\beta}
$$

\section*{Example}
\begin{itemize}
  \item Output grows at rate $1+g$.
  \item The interest rate sequence is constant over time:
\end{itemize}

$$
\tilde{R}(t)=\tilde{R}=\frac{1+g}{\beta}
$$

\begin{itemize}
  \item Why does the growth rate of output affects the interest rate?
\end{itemize}

\end{document}